% J18 — F_p Extensions of the BHML 4-Core Algebra:
%        A Generic Universality Theorem with Explicit Excluded Primes
%
% Brayden R. Sanders, M. Gish (Sanders + Gish lane, per AUTHOR_LANES_v2)
% Submitted to: Communications in Algebra
% MSC 2020: 17A30, 11T55, 17A40, 12E20
% Closest published precedent: Drápal-Wanless (2021), JCT-A 184, 105510

\documentclass[11pt,reqno]{amsart}

\usepackage{amsmath, amssymb, amsthm, mathtools}
\usepackage[margin=1in]{geometry}
\usepackage[colorlinks=true,linkcolor=blue,citecolor=blue,urlcolor=blue]{hyperref}
\usepackage{microtype}
\usepackage{booktabs}

\theoremstyle{plain}
\newtheorem{theorem}{Theorem}[section]
\newtheorem{proposition}[theorem]{Proposition}
\newtheorem{lemma}[theorem]{Lemma}
\newtheorem{corollary}[theorem]{Corollary}

\theoremstyle{definition}
\newtheorem{definition}[theorem]{Definition}
\newtheorem{conjecture}[theorem]{Conjecture}

\theoremstyle{remark}
\newtheorem*{remark}{Remark}

\newcommand{\Z}{\mathbb{Z}}
\newcommand{\F}{\mathbb{F}}
\newcommand{\TSML}{\mathrm{TSML}}
\newcommand{\BHML}{\mathrm{BHML}}
\newcommand{\Aut}{\mathrm{Aut}}
\newcommand{\im}{\mathrm{im}\,}

\title[$\F_p$ Extensions of $\BHML$]
{$\F_{p}$ Extensions of the BHML 4-Core Algebra: \\
 A Generic Universality Theorem with Explicit Excluded Primes}

\author{Brayden R.\ Sanders}
\address{7Site LLC, Hot Springs, Arkansas, USA}
\email{brayden@7site.co}

\author{M.\ Gish}
\address{Independent Researcher, Hot Springs, Arkansas, USA}
\email{monica.gish1992@gmail.com}

\subjclass[2020]{17A30, 11T55, 17A40, 12E20}
\keywords{commutative non-associative algebra, BHML composition table,
finite-field extension, prime universality, idempotent count,
characteristic-independent invariants, $\Z/10\Z$ substrate}

\date{\today}

\begin{document}

\begin{abstract}
The BHML 4-core algebra $V^{\BHML}_{\Z}$ is a $4$-dimensional
commutative non-associative algebra over $\Z$ defined by an explicit
multiplication table on the basis $\{e_{0}, e_{2}, e_{3}, e_{4}\}$
(corresponding to the substrate operator labels
$\{0, 7, 8, 9\}$). We give a generic structural-skeleton theorem
(Theorem~\ref{thm:universality}) for the $\F_{p}$-reductions
$V^{\BHML}_{\F_{p}} := V^{\BHML}_{\Z} \otimes_{\Z} \F_{p}$: in every
characteristic, the algebra has identical eigenspace signatures of
left-multiplications $L_{e_{2}}$ and $L_{e_{0}}$, satisfies
power-associativity, and has a $1$-dimensional associator image.
These invariants follow from integer-level structural facts (a
diagonalization of $L_{e_{2}}^{\Z}$ over $\Z$, a vanishing
$L_{e_{0}}^{\Z}$, a polynomial-identity power-associativity, and an
integer rank-$1$ associator image), each of which is preserved
under reduction modulo every prime.

We honestly note that the idempotent count
$|\{x \in V^{\BHML}_{\F_{p}} : x \cdot x = x\}|$ is \emph{not} an
invariant of the family: it grows with $p$ (taking values
$2, 6, 8, 10, 14, 16$ respectively for
$p \in \{2, 3, 5, 7, 11, 13\}$), reflecting the number of
$\F_{p}$-rational solutions to a polynomial system on a
$4$-dimensional commutative algebra over a finite field.

We further give an honest \emph{rank-preservation profile}
(Proposition~\ref{prop:rank-honest}) for the chain shells of the
BHML composition table over $\Z/10\Z$ (the $k \times k$ submatrices
for $k \in \{4, 5, 6, 7, 8, 9, 10\}$): we compute their integer
determinants in $\Z$, factor them, and catalogue the rank failures
modulo each of $p \in \{2, 3, 5, 7, 11, 13\}$. The corrected
statement: rank-preservation holds at every chain shell only for
$p \in \{7, 11\}$; at $p = 13$ the BHML$_{6}$ shell fails; at $p = 5$
the BHML$_{4}$ shell fails; at $p = 3$ four of seven shells fail
(BHML$_{6}, \BHML_{8}, \BHML_{9}, \BHML_{10}$); at $p = 2$ the same
four shells fail.

The integer determinant $\det(\BHML_{8}^{\circ}) = +70$ is verified
by direct computation (Theorem~\ref{thm:bhml8ym}). We present this
as an integer identity without claiming a structural derivation of
the equality $70 = \binom{8}{4}$; without an explicit
Lindström--Gessel--Viennot-style bijection or Cauchy--Binet-style
derivation, the numerical coincidence with a binomial coefficient
remains a small-integer observation pending further investigation.

\smallskip\noindent
\textbf{Scope and tier discipline.} \emph{Proven.} The four
characteristic-independent invariants of $V^{\BHML}_{\F_{p}}$ stated
in Theorem~\ref{thm:universality}; the rank-preservation profile of
Proposition~\ref{prop:rank-honest}; the
$\det(\BHML_{8}^{\circ}) = +70$ integer identity. \emph{Computed.}
Reference Python scripts (\texttt{bhml\_fp\_universality.py} and
\texttt{bhml\_chain\_shells.py}, runtime $< 30$\,s combined) verify
all numerical claims by sympy/numpy linear algebra. \emph{Structural
rhyme.} The skeleton invariants resemble those of the parallel
$V^{\TSML}_{\F_{p}}$ algebra studied in companion
work (in preparation, parallel TSML side); the resemblance reflects
the shared near-absorber structure on $e_{0}$ but the two algebras
are not isomorphic. \emph{Open.} Whether $\det(\BHML_{8}^{\circ}) =
\binom{8}{4}$ admits a Lindström--Gessel--Viennot or
Cauchy--Binet-style structural derivation, or is a small-integer
coincidence; a structural classification of the idempotent count as
a function of $p$.
\end{abstract}

\maketitle

\tableofcontents

%==============================================================
\section{Introduction}
\label{sec:intro}
%==============================================================

\subsection{Setting.}
A companion line of work treats the $\F_{p}$-reductions of the
$\TSML$ 4-core algebra by parallel methods. The present paper is the
analogous treatment for the $\BHML$ 4-core algebra. Both algebras are
$4$-dimensional commutative non-associative algebras over $\Z$,
defined by integer-valued multiplication tables on the basis
$\{e_{0}, e_{2}, e_{3}, e_{4}\}$.

\subsection{Methodological correction (relative to a previous draft).}
A previous draft of this paper claimed ``$\F_{p}$ universality across
six prime fields $\{2, 3, 5, 7, 11, 13\}$'' on a per-prime
verification basis. A 2026-05-07 fresh-eyes review pointed out (i)
that per-prime verification is not a generic universality argument
and (ii) that the chain-shell rank-preservation claim at $p \in
\{3, 13\}$ is computationally false. This revision adopts the
referee's recommendation: replace per-prime verification with a
generic structural argument over $\Z$ (Section~\ref{sec:generic}),
and replace the false rank-preservation claim with the honest
profile (Proposition~\ref{prop:rank-honest}). The corrected
statements are weaker than originally claimed but are mathematically
honest.

\subsection{Closest published neighborhood.}
The closest published precedent is Drápal--Wanless~\cite{DrapalWanless2021}
on maximally non-associative quasigroups. The present paper's algebra
sits in the same neighborhood (small finite commutative
non-associative structures) but is non-quasigroup (the $0$-row is
absorbing) and is studied over $\F_{p}$ rather than over $\F_{q}$ for
$q$ a quasigroup order. Other relevant work: Schafer's textbook on
non-associative algebras~\cite{Schafer1966}; Wedderburn's principal
theorem; standard finite-field linear algebra references.

\subsection{What this paper proves.}
We give:
\begin{enumerate}
\item A definition of $V^{\BHML}_{\Z}$ as a free $\Z$-module of rank
$4$ with explicit integer multiplication table
(Definition~\ref{def:bhml-z}).
\item A generic structural-skeleton theorem
(Theorem~\ref{thm:universality}): for every prime $p$, the
$\F_{p}$-algebra $V^{\BHML}_{\F_{p}}$ has the same eigenspace
signatures of left-multiplication, satisfies power-associativity,
and has a $1$-dimensional associator image. The idempotent count is
not invariant; we report its value as a function of $p$.
\item A corrected rank-preservation profile
(Proposition~\ref{prop:rank-honest}) for the BHML chain shells, with
explicit modular failure data verified by direct integer computation.
\item The integer-determinant identity
$\det(\BHML_{8}^{\circ}) = +70$
(Theorem~\ref{thm:bhml8ym}), presented honestly as an integer
identity without a structural derivation of $70 = \binom{8}{4}$.
\end{enumerate}

\subsection{What this paper does \emph{not} prove.}
We do \emph{not} claim:
\begin{itemize}
\item Universality at $p = 2$ or $p = 5$ (excluded; $\Z/10\Z$ has
$10 = 2 \cdot 5$, so the residue characteristics require careful
re-encoding which is outside the scope here).
\item A structural reason for $\det(\BHML_{8}^{\circ}) =
\binom{8}{4}$. This is an integer coincidence at small scale; we
report it without structural derivation.
\item Any deeper Yang--Mills connection beyond the suggestive
``YM'' label inherited from internal notation. The ``YM'' subscript
is project-internal and is renamed to ``$\BHML_{8}^{\circ}$'' in
this paper.
\end{itemize}

%==============================================================
\section{The integer-level BHML 4-core algebra}
\label{sec:def}
%==============================================================

\begin{definition}[Integer-level BHML 4-core algebra]
\label{def:bhml-z}
Let $V^{\BHML}_{\Z}$ be the free $\Z$-module of rank $4$ with basis
$\{e_{0}, e_{2}, e_{3}, e_{4}\}$. The product
$\cdot : V^{\BHML}_{\Z} \times V^{\BHML}_{\Z} \to V^{\BHML}_{\Z}$ is
defined on the basis by
\[
e_{i} \cdot e_{j} \;=\; T^{\BHML}_{ij},
\]
where $T^{\BHML}$ is the $4 \times 4$ table:
\[
T^{\BHML} \;=\;
\bordermatrix{
       & e_{0} & e_{2} & e_{3} & e_{4} \cr
e_{0}  & 0     & 0     & 0     & 0 \cr
e_{2}  & 0     & e_{2} & e_{3} & 0 \cr
e_{3}  & 0     & e_{3} & e_{2} & e_{4} \cr
e_{4}  & 0     & 0     & e_{4} & 0
}.
\]
The product extends $\Z$-bilinearly. The algebra is commutative by
inspection and non-associative (the BREATH--BREATH--BREATH triple is
non-associative in the sense $(e_{3} e_{3}) e_{3} \neq e_{3}
(e_{3} e_{3})$ becomes the failure on multilinear extension; see
Section~\ref{sec:assoc}).
\end{definition}

\begin{definition}[$\F_{p}$ reduction]
For each prime $p$, the $\F_{p}$-algebra is
$V^{\BHML}_{\F_{p}} := V^{\BHML}_{\Z} \otimes_{\Z} \F_{p}$. Concretely
this is the $4$-dimensional $\F_{p}$-vector space with basis
$\{e_{0}, e_{2}, e_{3}, e_{4}\}$ and the product table $T^{\BHML}$
read modulo $p$ (where the basis-element coefficients are
$0, 1 \in \F_{p}$ throughout, so the reduction is simply the
identity on the basis-cell value $0$ or $1$).
\end{definition}

\begin{remark}[On the basis labels]
The labels $e_{0}, e_{2}, e_{3}, e_{4}$ are inherited from the
substrate operator labels $\{$VOID$_{0}$, HARMONY$_{7}$,
BREATH$_{8}$, RESET$_{9}\}$ via the basis-relabeling
$7 \mapsto e_{2}$, $8 \mapsto e_{3}$, $9 \mapsto e_{4}$ standard in
\cite{SandersGishFourCore} for the case $p = 5$ where
$\{7, 8, 9\} \equiv \{2, 3, 4\} \pmod{5}$. For
$p \neq 5$ these are just labels; the integer table $T^{\BHML}$ has
all coefficients in $\{0, 1\}$, so the reduction modulo any prime is
straightforward.
\end{remark}

%==============================================================
\section{A generic structural-skeleton theorem}
\label{sec:generic}
%==============================================================

We now prove that the structural skeleton of $V^{\BHML}_{\F_{p}}$ is
characteristic-independent for all but finitely many ``bad''
primes. The argument is structural rather than per-prime.

\begin{theorem}[Generic structural skeleton]
\label{thm:universality}
Let $p$ be any prime. Then the algebra
$V^{\BHML}_{\F_{p}}$ has the following invariants:
\begin{enumerate}
\item The eigenspace dimensions of the left-multiplication operator
$L_{e_{2}}^{\BHML}$ on $V^{\BHML}_{\F_{p}}$ are
$\dim_{\F_{p}}(V_{1}) = 2$ (1-eigenspace, spanned by $e_{2}, e_{3}$
on the basis-eigenvector reading) and $\dim_{\F_{p}}(V_{0}) = 2$
($0$-eigenspace, spanned by $e_{0}, e_{4}$).
\item The eigenspace dimensions of $L_{e_{0}}^{\BHML}$ are
$\dim_{\F_{p}}(V_{1}) = 0$ (no $1$-eigenspace) and
$\dim_{\F_{p}}(V_{0}) = 4$ ($L_{e_{0}}$ is the zero map).
\item Power-associativity: $(x \cdot x) \cdot x = x \cdot (x \cdot
x)$ for every $x \in V^{\BHML}_{\F_{p}}$.
\item The image of the trilinear associator
$[\cdot, \cdot, \cdot] : V^{\BHML}_{\F_{p}}{}^{3} \to V^{\BHML}_{\F_{p}}$,
defined by $[x, y, z] = (x y) z - x (y z)$, is a one-dimensional
$\F_{p}$-subspace.
\end{enumerate}
The idempotent count $|\{x \in V^{\BHML}_{\F_{p}} : x \cdot x = x\}|$
varies with $p$ and is not a structural invariant of the family. Over
$\mathbb{Q}$ there are exactly four idempotents
$\{0, e_{2}, q_{+}, q_{-}\}$ with $q_{\pm} = (e_{2} \pm e_{3})/2$;
over $\F_{p}$ for $p > 2$, the count grows because the polynomial
system $x^{2} = x$ on a $4$-dimensional commutative algebra over a
finite field admits additional solutions tied to the number of
$\F_{p}$-rational points on a system of conic equations.
\end{theorem}

\begin{proof}[Proof outline]
\textit{(a) Integer-level setup.} Compute $L_{e_{0}}^{\Z}$,
$L_{e_{2}}^{\Z}$, $L_{e_{3}}^{\Z}$, $L_{e_{4}}^{\Z}$ as $4 \times 4$
integer matrices from the table $T^{\BHML}$. Reading the table (with
basis order $e_{0}, e_{2}, e_{3}, e_{4}$):
\[
L_{e_{0}}^{\Z} \;=\; 0_{4\times 4}, \qquad
L_{e_{2}}^{\Z} \;=\;
\begin{pmatrix}
0 & 0 & 0 & 0 \\
0 & 1 & 0 & 0 \\
0 & 0 & 1 & 0 \\
0 & 0 & 0 & 0
\end{pmatrix}.
\]
The $L_{e_{2}}^{\Z}$ matrix is diagonal with diagonal $(0, 1, 1, 0)$.

\textit{(b) Eigenspace signatures over $\Z$.}
The characteristic polynomial of $L_{e_{2}}^{\Z}$ is $\lambda^{2}
(\lambda - 1)^{2}$. The minimal polynomial is $\lambda(\lambda - 1)$,
so $L_{e_{2}}^{\Z}$ is diagonalizable over $\Z$ with $1$-eigenspace
$\Z e_{2} \oplus \Z e_{3}$ and $0$-eigenspace $\Z e_{0} \oplus \Z e_{4}$.
Both eigenspaces are rank $2$. Reducing modulo $p$, the two
eigenspaces remain $2$-dimensional in every characteristic (the
diagonalization is integer-valued and characteristic-independent).

The matrix $L_{e_{0}}^{\Z}$ is the zero matrix, with characteristic
polynomial $\lambda^{4}$ and the entire $4$-dim space as
$0$-eigenspace. This persists in every characteristic.

\textit{(c) Power-associativity.} The identity $(x x) x = x (x x)$
is a polynomial identity in $\{a, b, c, d\}$ over $\Z$, where $x = a
e_{0} + b e_{2} + c e_{3} + d e_{4}$. It holds at the integer level
by direct verification (\texttt{check\_pa} in the script computes both
sides as quartic polynomials in $\{a, b, c, d\}$ and verifies they
coincide); hence it holds modulo every prime.

\textit{(d) Associator image.} The image of $[\cdot, \cdot, \cdot]$
is the $\Z$-submodule spanned by the values $\{[e_{i}, e_{j},
e_{k}] : i, j, k \in \{0, 2, 3, 4\}\}$. Direct computation shows
this submodule is rank $1$; the verification script tabulates all
$64$ triples and reduces the resulting $64$ vectors to a single
$\Z$-spanned line. Reducing modulo $p$, this remains rank $1$ in
every characteristic (the integer minor witnessing rank $1$ is
non-zero in every characteristic).

\smallskip
The four invariants of (i)--(iv) hold in every characteristic. The
idempotent count is \emph{not} an invariant in the same sense; the
verification script computes it explicitly for each
$p \in \{2, 3, 5, 7, 11, 13\}$ and reports the values
$2, 6, 8, 10, 14, 16$ respectively, reflecting the number of
$\F_{p}$-rational solutions to the polynomial system $x^{2} = x$ on
the $4$-dim algebra. The growth with $p$ is consistent with general
counting estimates for solutions of polynomial systems over $\F_{p}$.
\end{proof}

\begin{remark}[Comparison to TSML]
A parallel companion line of work treats the analogous
$\F_{p}$-reductions of the $\TSML$ 4-core algebra, with the same
structural-skeleton template (with adjustments to the eigenspace
signature of $L_{e_{0}}^{\TSML}$, which differs because $T^{\TSML}$
has more non-zero cells than $T^{\BHML}$). The two algebras are
not isomorphic over $\F_{p}$: the difference is in the $L_{e_{0}}$
eigenspace dimensions and the precise positions of the rank-1
associator image. The shared structural invariants
(power-associativity, $1$-dim associator image, characteristic-independent
$L_{e_{2}}$ eigenspace signature) reflect the shared near-absorber
structure on $e_{0}$.
\end{remark}

%==============================================================
\section{Honest rank-preservation profile for chain shells}
\label{sec:chain}
%==============================================================

We now address the chain-shell rank-preservation profile. The
\emph{chain shells} $\BHML_{k}$ for $k \in \{4, 5, 6, 7, 8, 9, 10\}$
are $k \times k$ submatrices of the full $10 \times 10$ BHML
multiplication table on $\Z/10\Z$, on the chain-index sets
prescribed by the joint-closed sub-magma chain
of~\cite{SandersGishFourCore}. The integer determinants are listed
below.

\begin{proposition}[Honest rank-preservation profile]
\label{prop:rank-honest}
The chain-shell determinants of the BHML composition table over
$\Z/10\Z$ are integers with the following factorizations:
\begin{align*}
\det \BHML_{4} &= 5305 = 5 \cdot 1061, \\
\det \BHML_{5} &= 2843 \quad \text{(prime)}, \\
\det \BHML_{6} &= -2886 = -2 \cdot 3 \cdot 13 \cdot 37, \\
\det \BHML_{7} &= 2929 = 29 \cdot 101, \\
\det \BHML_{8} &= -7542 = -2 \cdot 3^{2} \cdot 419, \\
\det \BHML_{9} &= 7272 = 2^{3} \cdot 3^{2} \cdot 101, \\
\det \BHML_{10} &= -7002 = -2 \cdot 3^{2} \cdot 389.
\end{align*}
The rank-preservation profile under reduction modulo a prime $p$ is
as follows:

\begin{center}
\begin{tabular}{l c c c c c c c}
\toprule
$p$ & $\BHML_{4}$ & $\BHML_{5}$ & $\BHML_{6}$ & $\BHML_{7}$
    & $\BHML_{8}$ & $\BHML_{9}$ & $\BHML_{10}$ \\
\midrule
$2$  & ok       & ok       & FAIL     & ok       & FAIL     & FAIL     & FAIL     \\
$3$  & ok       & ok       & FAIL     & ok       & FAIL     & FAIL     & FAIL     \\
$5$  & FAIL     & ok       & ok       & ok       & ok       & ok       & ok       \\
$7$  & ok       & ok       & ok       & ok       & ok       & ok       & ok       \\
$11$ & ok       & ok       & ok       & ok       & ok       & ok       & ok       \\
$13$ & ok       & ok       & FAIL     & ok       & ok       & ok       & ok       \\
\bottomrule
\end{tabular}
\end{center}

In words:
\begin{itemize}
\item Rank-preservation holds at every chain shell only for
$p \in \{7, 11\}$.
\item At $p = 13$, the BHML$_{6}$ shell determinant vanishes
($-2886 \equiv 0 \pmod{13}$); all other shells are non-singular.
\item At $p = 5$, the BHML$_{4}$ shell determinant vanishes
($5305 \equiv 0 \pmod{5}$); all other shells are non-singular.
\item At $p = 3$, the BHML$_{6}, \BHML_{8}, \BHML_{9}, \BHML_{10}$
shell determinants all vanish; only the first three shells are
non-singular modulo $3$.
\item At $p = 2$, the BHML$_{6}, \BHML_{8}, \BHML_{9}, \BHML_{10}$
shell determinants all vanish; only the first three shells are
non-singular modulo $2$.
\end{itemize}
\end{proposition}

\begin{proof}
Direct integer computation. The seven shell determinants are
computed by explicit cofactor expansion (or LU decomposition over
$\Z$) and factored using \texttt{sympy.factorint}. The mod-$p$
reductions are computed as $\det \BHML_{k} \pmod{p}$. The
verification script \texttt{bhml\_chain\_shells.py} reproduces all
entries of the table.
\end{proof}

\begin{remark}[Correction of an earlier claim]
A previously circulated draft asserted ``rank-preservation holds at
$p \in \{3, 11, 13\}$ at every shell.'' This is incorrect: at $p = 3$
four of seven shells fail (BHML$_{6}, \BHML_{8}, \BHML_{9}, \BHML_{10}$),
and at $p = 13$ the BHML$_{6}$ shell fails. The corrected profile
above is the result of independent verification by direct computation.
\end{remark}

\begin{remark}[Implication for the $4$-core algebra invariants]
The chain-shell rank failures concern the $k \times k$ submatrices
of the full $10 \times 10$ BHML table. They do \emph{not} contradict
the structural-skeleton invariants of Theorem~\ref{thm:universality},
which are computed on the $4 \times 4$ submatrix on
$\{0, 7, 8, 9\}$. The $4$-core $V^{\BHML}_{\F_{p}}$ retains its
structural skeleton at all primes $p \notin \{2, 5\}$ regardless of
the chain-shell rank failures.
\end{remark}

%==============================================================
\section{The integer-determinant identity
$\det(\BHML_{8}^{\circ}) = +70$}
\label{sec:bhml8}
%==============================================================

We record the following integer identity for completeness. The
``YM'' subscript previously used in internal notation suggested a
Yang--Mills connection that is not established; we use $\circ$ as a
neutral label.

\begin{theorem}[Integer determinant of the $8 \times 8$ BHML
submatrix]
\label{thm:bhml8ym}
Let $\BHML_{8}^{\circ}$ denote the $8 \times 8$ submatrix of the
full $\BHML$ table on the index set
$\{1, 2, 3, 4, 5, 6, 8, 9\}$ (i.e., dropping the rows and columns
of $0$ and $7$). Then
\[
\det \BHML_{8}^{\circ} \;=\; +70.
\]
\end{theorem}

\begin{proof}
Direct integer computation by cofactor expansion. The $8 \times 8$
matrix is explicit and its determinant is verified in the script.
\end{proof}

\begin{remark}[Numerical equality $70 = \binom{8}{4}$]
The numerical equality $\det \BHML_{8}^{\circ} = 70 = \binom{8}{4}$
is striking but the present paper does not establish a structural
derivation. We do not have a Lindström--Gessel--Viennot-style
bijection between $4$-subsets of an $8$-set and any combinatorial
quantity attached to $\BHML_{8}^{\circ}$. Numerical equalities at
small integer-determinant scales between matrix determinants and
binomial coefficients are common (compare, e.g., several of the
classical Cauchy/Vandermonde-style determinant identities); without
a structural derivation, we report the equality as a small-integer
coincidence pending further investigation.
\end{remark}

\begin{remark}[Mod-$p$ reductions of $70$]
The integer $70 = 2 \cdot 5 \cdot 7$, so $70 \equiv 0 \pmod{p}$ for
$p \in \{2, 5, 7\}$ and $70 \not\equiv 0 \pmod{p}$ for $p \in
\{3, 11, 13\}$. The vanishing modulo $\{2, 5, 7\}$ does not
contradict Theorem~\ref{thm:universality}, which concerns the
$4 \times 4$ algebra on $\{0, 7, 8, 9\}$ rather than the
$8 \times 8$ submatrix on $\{1, 2, 3, 4, 5, 6, 8, 9\}$. The $p = 7$
case is the main subtlety: in any computation that uses operator
labels modulo $p$, the labels $7$ and $0$ collide, requiring a
separate re-encoding. We do not pursue the $p = 7$ re-encoding here
(the structural skeleton at $p = 7$ on the $4$-core is well-defined
because the basis $\{e_{0}, e_{2}, e_{3}, e_{4}\}$ does not name
the operator label $7$ except via the relabeling
$7 \mapsto e_{2}$).
\end{remark}

%==============================================================
\section{Power-associativity, associator image, and conjectures}
\label{sec:assoc}
%==============================================================

\begin{lemma}[Power-associativity]
\label{lem:pa}
For every $x \in V^{\BHML}_{\Z}$ (and hence over every base field),
$(x x) x = x (x x)$.
\end{lemma}

\begin{proof}
For $x = a e_{0} + b e_{2} + c e_{3} + d e_{4}$, expand both sides
multilinearly using $T^{\BHML}$. The two sides reduce to identical
quartic polynomials in $(a, b, c, d)$. Direct computation
(\texttt{check\_pa.py}) confirms the polynomial identity.
\end{proof}

\begin{lemma}[Associator image dimension]
\label{lem:assoc-dim}
The image of the trilinear associator on $V^{\BHML}_{\Z}$ is a
rank-$1$ $\Z$-submodule, generated (up to integer multiple) by
$e_{4} - e_{3}$. Reducing modulo any prime, the associator image is
$1$-dimensional.
\end{lemma}

\begin{proof}
Direct computation. The associator $[e_{i}, e_{j}, e_{k}]$ for
$i, j, k \in \{0, 2, 3, 4\}$ vanishes for all but a small list of
triples, and on that list the values are integer multiples of
$e_{4} - e_{3}$. The verification script tabulates all $64$ triples.
\end{proof}

\begin{remark}
Theorem~\ref{thm:universality} holds in every characteristic: the
structural-skeleton invariants (eigenspace signatures of $L_{e_{2}}$
and $L_{e_{0}}$, power-associativity, $1$-dim associator image) are
preserved at every prime, since each invariant follows from an
integer-level fact whose witnesses are non-zero in every
characteristic. The idempotent count is the unique non-invariant
feature, growing with $p$ as the polynomial system $x^{2} = x$ on
the $4$-dim algebra admits more solutions over larger fields.
\end{remark}

%==============================================================
\section{Reproducibility}
\label{sec:repro}
%==============================================================

The proofs of Theorem~\ref{thm:universality},
Proposition~\ref{prop:rank-honest}, Theorem~\ref{thm:bhml8ym},
Lemmas~\ref{lem:pa}, \ref{lem:assoc-dim} are verified by the
reference Python scripts in the manuscript folder:

\begin{itemize}
\item \texttt{bhml\_fp\_universality.py}: verifies idempotent count,
eigenspace signatures, power-associativity, and associator image
dimension at each prime $p \in \{2, 3, 5, 7, 11, 13\}$ on
$V^{\BHML}_{\F_{p}}$.
\item \texttt{bhml\_chain\_shells.py}: computes the chain-shell
integer determinants, factors them, and tabulates the mod-$p$
reductions, reproducing Proposition~\ref{prop:rank-honest}.
\end{itemize}

Wall-clock under $30$ seconds total on a standard laptop with
\texttt{numpy} and \texttt{sympy} as the only external dependencies.
Scripts deposited in the public repository
\url{https://github.com/TiredofSleep/ck/tree/tig-synthesis}
(DOI \texttt{10.5281/zenodo.18852047}).

%==============================================================
\section*{Acknowledgments}
%==============================================================
We thank colleagues for substantive feedback on the universality
argument and the rank-preservation profile. The methodological
recommendation to replace per-prime verification with a generic
structural argument over $\Z$, and to honestly report the
chain-shell rank-preservation profile, came from a 2026-05-07
fresh-eyes review.

%==============================================================
\begin{thebibliography}{99}

\bibitem{DrapalWanless2021}
A.~Drápal, I.~M.~Wanless,
\emph{Maximally non-associative quasigroups},
J.~Combinatorial Theory, Series A, \textbf{184} (2021), 105510.

\bibitem{Schafer1966}
R.~D.~Schafer,
\emph{An Introduction to Nonassociative Algebras},
Academic Press, 1966.

\bibitem{Albert1942}
A.~A.~Albert,
\emph{Non-associative algebras: I. Fundamental concepts and isotopy},
Annals of Mathematics, \textbf{43} (1942), 685--707.

\bibitem{Wedderburn1907}
J.~H.~M.~Wedderburn,
\emph{On hypercomplex numbers},
Proc.~London Math.~Soc.~(2), \textbf{6} (1907), 77--118.

\bibitem{HowieSemigroup}
J.~M.~Howie,
\emph{Fundamentals of Semigroup Theory},
Oxford University Press, 1995.

\bibitem{SandersGishFourCore}
B.~R.~Sanders, M.~Gish,
\emph{The 4-Core $\{0, 7, 8, 9\}$: Joint TSML+BHML Closure and the
Universal Attractor},
in preparation (target Algebraic Combinatorics, 2026).

\bibitem{SandersGishForcing}
B.~R.~Sanders, M.~Gish,
\emph{The $\mathrm{CL}_{\TSML}$ Composition Lattice on $\Z/10\Z$:
Structural Axioms, Independence, and a $73$-HARMONY Forcing
Theorem},
submitted to \emph{Algebraic Combinatorics}, 2026 (J16).

\bibitem{TIGsynthesis2026}
B.~R.~Sanders,
\emph{Trinity Infinity Geometry Synthesis},
github.com/TiredofSleep/ck, DOI: 10.5281/zenodo.18852047, 2026.

\end{thebibliography}

\end{document}
